\chapter{Analysis of the subject area and problem statement}\label{chapter-1.-analysis-of-the-subject-area-and-problem-statement}

This chapter establishes the technical and scientific context of the work. Section 1.1 describes urban-tree inventory as an applied problem. Section 1.2 surveys the remote-sensing data sources available for the task. Section 1.3 organises the relevant deep-learning literature of 2019 -- 2025 into a single comparative table grouped by model family. Section 1.4 identifies the \textbf{geographic generalisation gap} as the dominant obstacle to applying off-the-shelf models in Astana. Section 1.5 formulates the precise problem statement that the rest of the work will solve.

\section{Urban-tree inventory as an applied problem}\label{urban-tree-inventory-as-an-applied-problem}

Municipalities increasingly rely on quantitative inventories of urban green space for planning, environmental monitoring and climate adaptation. The basic unit of such inventories is the individual tree, characterised by a geographic position, a crown size and --- when available --- a species label and a health status. Aggregated indicators (number of trees per district, canopy-cover percentage, average crown area) follow directly from the per-tree records.

Historically, urban-tree inventories were obtained through manual field surveys with hand-held GPS receivers and paper forms. Reports from comparable cities --- for example, the multi-year inventory of Sofia, Bulgaria {[}6{]} or the multi-month effort of Pasadena, California {[}18{]} --- confirm that for any large city the manual approach is slow, expensive and rapidly becomes obsolete. Three classes of automated alternative have emerged: airborne LiDAR {[}23{]}, which is the gold standard wherever it is available but is exceptionally expensive and unsuitable for cities --- including Astana --- that do not appear in national LiDAR programmes; high-resolution aerial or UAV imagery at 5 -- 30 cm GSD, the dominant data source for urban-tree research {[}10, 11, 4{]}; and \textbf{freely or inexpensively available very-high-resolution satellite imagery} from commercial vendors and aggregators such as ESRI World Imagery and Google Earth --- the only practical option for Astana at the required scale {[}14, 5, 2{]}.

The practical requirements imposed by \emph{Zelenstroy} on any automated inventory are: (i) accept imagery from heterogeneous sources without re-training; (ii) output a per-tree polygon mask rather than a single pixel, so that crown area can be computed; (iii) convert pixel coordinates into geographic coordinates in WGS-84; (iv) export the inventory in formats consumable by QGIS, ArcGIS and Excel; (v) provide a confidence score for each detection so that human reviewers can prioritise their inspection time. The methodological choices documented in Chapter 2 are direct consequences of these five requirements.

\section{Remote-sensing data sources}\label{remote-sensing-data-sources}

The deep-learning literature on tree detection is highly \textbf{data-source-dependent}: a model that excels on 5-cm UAV imagery typically fails when applied directly to 1-m satellite imagery. The body of work surveyed below is organised along this dimension in Table 1.1.

\textbf{Table 1.1 --- Data sources reported in the surveyed literature.}

{\def\LTcaptype{none} % do not increment counter
\begin{longtable}[]{@{}>{\raggedright\arraybackslash}p{0.237\linewidth} >{\raggedright\arraybackslash}p{0.139\linewidth} >{\raggedright\arraybackslash}p{0.287\linewidth} >{\raggedright\arraybackslash}p{0.109\linewidth} >{\raggedright\arraybackslash}p{0.178\linewidth}@{}}
\toprule\noalign{}
\textbf{Source} & \textbf{Typical GSD} & \textbf{Coverage} & \textbf{Cost} & \textbf{Representative works} \\
\midrule\noalign{}
\endhead
\bottomrule\noalign{}
\endlastfoot
Airborne LiDAR & 5 -- 50 cm 3-D & National / regional campaigns & Very high & {[}23, 1{]} \\
UAV / drone RGB & 1 -- 5 cm & A single survey area & Medium & {[}11, 8, 13{]} \\
Aerial orthophotos & 5 -- 30 cm & Country-wide & Medium-high & {[}10, 4, 3{]} \\
Very-high-res satellite & 30 cm -- 1 m & Global & Low--medium & {[}2, 5, 14{]} \\
Sentinel-2 multispectral & 10 m & Global, free & Free & {[}14, 16, 17{]} \\
\end{longtable}
}

The very-high-resolution satellite category is the most recent in the literature --- most dedicated satellite papers appeared in 2024 and 2025 --- and is the only data source practically available for Astana.

\section{Deep-learning paradigms for tree detection}\label{deep-learning-paradigms-for-tree-detection}

Four model families are relevant to the present work: two-stage detectors (Faster R-CNN, Mask R-CNN), one-stage detectors (YOLO, RetinaNet), the domain-specialised DeepForest detector, and the SAM 2 foundation model for mask refinement. Table 1.2 consolidates the best reported numbers from each family.

\textbf{Two-stage detectors.} Faster R-CNN {[}27{]} and its instance-segmentation extension Mask R-CNN {[}28{]} were the first deep architectures applied to individual tree detection. The most recent contribution in this family is the MCAN architecture of Lv et al.~{[}8{]}, which adds a CSPNet backbone and CBAM attention to Mask R-CNN and reports a detection AP of 92.40 \% on UAV imagery of a forested campus in Zhejiang, China. The principal disadvantage of the family is \textbf{computational cost}: for an interactive web tool, Mask-R-CNN-class instance segmenters are too slow on a laptop GPU. This consideration motivates the choice of a one-stage architecture for the YOLO branch of the present system.

\textbf{One-stage detectors: YOLO and RetinaNet.} The YOLO family {[}30, 32{]} has been the dominant choice for urban-tree detection since 2022. Velasquez-Camacho et al.~{[}4{]} tested YOLOv5 against DeepForest and Faster R-CNN on aerial and satellite imagery of Lleida, Spain, and reported an F-score of 84.9 \% for the YOLOv5x variant --- substantially better than Faster R-CNN's 35.6 \%. The most recent benchmark of Abbas and Damaševičius {[}2{]} evaluated every modern YOLO variant --- YOLOv8 through YOLOv12 --- on a public RGB satellite tree dataset of 3 157 images and reported YOLOv12m as the best performer with mAP@50 = 90.8 \%. Sun {[}7{]} dedicated an entire PhD thesis to YOLOv7/v8 for tree-crown instance segmentation on Wellington, New Zealand aerial imagery, and reported that her improved YOLOv8 variants outperform Mask R-CNN, YOLOv5 and SOLOv2 on both Box AP and Mask AP with fewer parameters. The other relevant one-stage architecture is \textbf{RetinaNet} {[}29{]}: the benchmark by dos Santos et al.~{[}11{]} reported a RetinaNet AP of 92.64 \%, decisively outperforming YOLOv3 (85.88 \%) and Faster R-CNN (82.48 \%) on UAV imagery of Campo Grande, Brazil. This empirical result motivated the developers of DeepForest {[}1{]} to choose RetinaNet as their base architecture.

\textbf{The DeepForest model.} DeepForest {[}1{]} is a tree-specialised RetinaNet detector packaged as a Python library with pre-trained weights and a stable \texttt{predict\_tile()} API. The original release was trained on a large semi-supervised dataset derived from the NEON aerial-lidar campaigns over forested sites in the United States, and reported a baseline F1 ≈ 0.65 on held-out NEON scenes. Two empirical results from the literature decisively constrain its use on urban imagery. Ventura et al.~{[}3{]} tested DeepForest off-the-shelf on 60-cm NAIP imagery of eight Californian cities and reported precision 0.735 but \textbf{recall of only 0.294}, for an F-score of 0.42; after fine-tuning on a few hundred urban tiles, the same model recovered an F-score of \textbf{0.729}. The Sofia DeepForest work of Dakov and Petrova-Antonova {[}6{]} provides a complementary data point: trained on 826 manually-annotated trees in Sofia, Bulgaria, DeepForest reached an F1 of 0.674 -- 0.685 --- comparable to the original NEON benchmark but obtained on what is geographically and architecturally the closest analogue to Astana in the entire surveyed literature. Both works conclude that \textbf{DeepForest must be fine-tuned for urban use}.

\textbf{Foundation models: SAM and SAM 2.} The newest paradigm in the field is the foundation-model approach exemplified by the Segment Anything Model (SAM) {[}25{]} and its 2024 successor SAM 2 {[}36{]}. Both expose a prompt-based interface in which the user supplies a point, a bounding box or a coarse mask and the model returns a precise segmentation of the corresponding object. The critical property is \textbf{zero-shot generalisation}: the models produce sharp masks even on object categories that were never explicitly labelled at training time, including trees in satellite imagery. SAM 2 specifically can therefore be used as a mask-refinement stage on top of any bounding-box detector --- DeepForest in our case --- without requiring an additional fine-tune. This idea is the basis of the fourth model branch in Section 2.6.

\textbf{Table 1.2 --- Best reported results from the surveyed literature, sorted by method family.}

{\def\LTcaptype{none} % do not increment counter
\begin{longtable}[]{@{}>{\raggedright\arraybackslash}p{0.214\linewidth} >{\raggedright\arraybackslash}p{0.214\linewidth} >{\raggedright\arraybackslash}p{0.093\linewidth} >{\raggedright\arraybackslash}p{0.214\linewidth} >{\raggedright\arraybackslash}p{0.214\linewidth}@{}}
\toprule\noalign{}
\textbf{Method} & \textbf{Work} & \textbf{Year} & \textbf{Data} & \textbf{Best metric} \\
\midrule\noalign{}
\endhead
\bottomrule\noalign{}
\endlastfoot
Faster R-CNN & dos Santos et al.~{[}11{]} & 2019 & UAV RGB, Campo Grande & AP = 82.48 \% \\
Mask R-CNN (MCAN) & Lv et al.~{[}8{]} & 2023 & UAV RGB, Zhejiang & Det AP = 92.40 \%, Seg AP = 97.70 \% \\
RetinaNet & dos Santos et al.~{[}11{]} & 2019 & UAV RGB, Campo Grande & AP = 92.64 \% \\
YOLOv4-Lite & Zheng and Wu {[}12{]} & 2022 & Google Earth 0.27 m & Acc = 96.3 \% (campus) \\
YOLOv5x & Velasquez-Camacho et al.~{[}4{]} & 2023 & Aerial + sat., Lleida & F1 = 84.9 \% \\
YOLOv8 (Wellington) & Sun {[}7{]} & 2025 & Aerial, Wellington & Best Box and Mask AP vs.~Mask R-CNN / YOLOv5 / SOLOv2 \\
YOLOv12m & Abbas and Damaševičius {[}2{]} & 2025 & RGB satellite, 3 157 imgs & \textbf{mAP@50 = 90.8 \%}, mAP@50:95 = 58.1 \% \\
U-Net & Wang et al.~{[}9{]} & 2021 & Aerial 32 cm, Vaihingen & OA = 99.14 \%, IoU = 96.38 \% \\
DeepLabV3+ & Martins et al.~{[}10{]} & 2021 & Aerial 10 cm, Campo Grande & F1 = 91.4 \%, IoU = 73.89 \% \\
DeepForest off-the-shelf (urban) & Ventura et al.~{[}3{]} & 2024 & NAIP 60 cm, 8 CA cities & F = \textbf{0.42} (P = 0.74, R = 0.29) \\
DeepForest fine-tuned (urban) & Ventura et al.~{[}3{]} & 2024 & NAIP 60 cm, 8 CA cities & F = \textbf{0.729} \\
DeepForest urban (Sofia) & Dakov and Petrova-Antonova {[}6{]} & 2024 & Aerial 10 cm, Sofia & F1 = 0.674 -- 0.685 \\
Sub-pixel canopy (CASNet) & He et al.~{[}14{]} & 2022 & Sentinel-2, 34 Chinese cities & OA = 88.6 \% \\
Canopy height (MUFCH) & Xu et al.~{[}16{]} & 2025 & Sentinel-2 + OSM, Beijing & MAE = 2.02 m \\
\textbf{YOLOv8x-seg v4 (this work)} & this thesis & 2026 & Very-high-res sat., Astana M14 & \textbf{Box mAP@50 = 0.315, Mask mAP@50 = 0.290} \\
\textbf{Mask R-CNN v2+v3 (this work)} & this thesis & 2026 & Very-high-res sat., Astana M14 & Box mAP@50 = 0.166, Mask mAP@50 = 0.158 \\
\textbf{DeepForest v3 + SAM 2 (this work)} & this thesis & 2026 & Very-high-res sat., Astana M14 & Box mAP@50 = 0.146, Mask mAP@50 = 0.134 \\
\textbf{NEON off-the-shelf on Astana (this work)} & this thesis & 2026 & Very-high-res sat., Astana M14 & Box mAP@50 = \textbf{0.012} (first measurement on a Central-Asian city) \\
\end{longtable}
}

Three observations follow from this table. First, the best reported numbers in the literature are very high --- state-of-the-art detection models routinely exceed 90 \% mAP at IoU = 0.5. Second, \textbf{all the best results are obtained on datasets from the United States, Europe, China, Brazil or New Zealand}: there is no paper in the surveyed corpus that evaluates any of these methods on Central-Asian imagery, no paper that uses imagery of Kazakhstan, and no paper that addresses the specific architectural context of a Soviet-era micro-district city. Third, the gap between off-the-shelf and fine-tuned DeepForest is enormous --- 0.42 → 0.73 F1 on the same data after a few hundred annotations {[}3{]}. This is the single most important quantitative finding for the present work, because it directly justifies the time investment in building a custom Astana annotated dataset.

\section{The geographic generalisation gap}\label{the-geographic-generalisation-gap}

The pattern --- high in-domain performance, large performance drop on out-of-domain imagery --- is sometimes called the \textbf{geographic generalisation gap}. It is not specific to tree detection: similar drops have been reported in building segmentation, land-cover classification and road extraction whenever a model trained in one country is applied without fine-tuning to imagery of another.

For Astana the magnitude of the gap was unknown at the start of the project, since no published number existed for any deep-learning tree-detection model on Central-Asian satellite imagery. As part of this work the public DeepForest checkpoint (trained on NEON aerial-lidar campaigns of forested sites in the United States) was evaluated out-of-the-box on the Astana M14 validation set described in Section 3.1. The result is a Box mAP@50 of \textbf{0.012} --- one to two orders of magnitude below any of the fine-tuned configurations reported later, and three orders of magnitude below the same model's off-the-shelf 0.42 F-score on NAIP USA imagery {[}3{]}. This single number is the empirical anchor of the entire diploma project: it quantifies, for the first time on a Central-Asian capital, the gap that motivates the dataset-construction and fine-tuning effort of Chapters 2 and 3.

Four factors are commonly invoked to explain the gap: \textbf{floristic differences} (Astana's urban canopy is dominated by \emph{Populus}, elms, birches and apricots, which present visually different crowns from the broadleaf and conifer species typical of North-American NEON sites); \textbf{urban-morphology differences} (Soviet-era micro-district planning produces a regular grid of multi-storey apartment blocks separated by narrow strips of green space, qualitatively different from suburban-sprawl U.S. benchmarks and dense historic-city European benchmarks); \textbf{resolution and acquisition-geometry mismatch} (satellite look-angles can deviate by up to 30° from nadir, leading to large variations in tree-shadow direction within a single image); and \textbf{annotation-policy differences} (per-crown vs.~per-cluster labelling conventions vary between datasets). The conclusion is that off-the-shelf models cannot be expected to deliver state-of-the-art numbers on Astana imagery and that some form of domain adaptation --- fine-tuning, ensemble or both --- is mandatory. This is the principal motivation for the methodology of Chapter 2.

\section{Problem statement}\label{problem-statement}

Based on the analysis above, the problem solved in the present work can be stated formally as follows.

\textbf{Given:} a colour satellite image of an arbitrary area of Astana at a ground sampling distance of approximately 0.3 -- 1 m per pixel, optionally accompanied by geographic metadata (a GeoTIFF affine transform, two or four corner coordinates supplied by the user, or none).

\textbf{Produce:} an inventory \(I = \{(p_i, c_i, s_i, a_i, l_i, \lambda_i, \phi_i)\}_{i=1}^{N}\) of \(N\) trees, where for each tree \(i\), \(p_i\) is a polygon mask approximating the projected crown, \(c_i\) is the bounding box of the polygon in pixel coordinates, \(s_i \in [0, 1]\) is a confidence score, \(a_i\) is the crown area in pixels and (when geographic conversion is possible) in square metres, \(l_i\) is the class label (restricted to the single class ``Tree''), and \((\lambda_i, \phi_i)\) are the longitude and latitude of the crown centroid in WGS-84 (defined when geographic conversion is possible).

\textbf{Subject to:} per-image inference time of at most 30 seconds on a single laptop GPU of the GeForce RTX 4060 class with 8 GiB of VRAM; Box mAP@50 on a held-out Astana validation set strictly above the public NEON-pretrained baseline of 0.012 (Section 1.4) by at least one order of magnitude, with the explicit aim of being the strongest published number on Astana satellite imagery to date; support for sliding-window tiled inference so that arbitrarily large input images can be processed; support for at least three export formats --- GeoJSON, CSV and standalone HTML --- to satisfy the interoperability requirement; and compatibility with arbitrary geographic input modes so that the system can be used both with rigorous GeoTIFF deliverables and with informal screenshots supplied by a non-technical user.

Chapter 2 presents the system design that meets these requirements, and Chapter 3 reports its experimental evaluation against the literature baselines of Table 1.2.

\newpage
